\section{Method}
\label{sec:method}

\subsection{Problem formulation and behavior cloning}
A demonstration $\tau_i=\{(s_{it},a_{it})\}_{t=1}^{T_i}$ contains observed
robot/task state $s$ and expert action $a=(a^{\mathrm{mot}},a^{\mathrm{grip}})$.
The motion component is a three-dimensional Cartesian displacement; the
gripper component is binary. A shared MLP produces a motion vector and one
gripper logit. With training-only observation statistics $\mu,\sigma$ and
motion scale $s_a$, we optimize
\begin{equation}
\mathcal{L}=
\left\|\hat a^{\mathrm{mot}}/s_a-a^{\mathrm{mot}}/s_a\right\|_2^2
+\operatorname{BCE}(\hat a^{\mathrm{grip}},a^{\mathrm{grip}}).
\end{equation}
At inference, motion is rescaled and norm-clipped to the largest training
expert step. Evaluation is autonomous and closed loop; expert actions are
never used during policy rollouts.

\subsection{Difficulty definitions}
All scores are trajectory-level and use only the 240 training trajectories.
For a component $x$, $z(x)=(x-\mu_x)/(\sigma_x+\epsilon)$ uses statistics from
that split.

\paragraph{Task difficulty.}
For pick-and-place, the components are initial end-effector--object distance,
object--goal distance, and inverse target radius $1/r$:
\begin{equation}
d_{\mathrm{task}}(\tau)=\tfrac13[z(d_{\mathrm{EE,obj}})
+z(d_{\mathrm{obj,goal}})+z(1/r)].
\end{equation}
They represent approach distance, transport distance, and placement precision.

\paragraph{Expert difficulty.}
We average standardized duration $T$, executed end-effector path length $L$,
and cumulative Cartesian action magnitude $A$:
\begin{equation}
d_{\mathrm{expert}}(\tau)=\tfrac13[z(T)+z(L)+z(A)].
\end{equation}
These strongly correlated components measure execution effort; $L$ and $A$
are especially redundant in our deterministic controller.

\paragraph{Learner difficulty.}
We partition the training trajectories into five folds. For each fold, a new
BC model with the same $128$--$128$ architecture trains for 95 fixed epochs on
192 trajectories with fold-only normalization and motion scaling, then scores
the unseen 48. This differs from the 20,600-update curriculum budget. Every
trajectory is scored exactly once:
\begin{equation}
d_{\mathrm{learner}}(\tau)=\frac{1}{T}\sum_{t=1}^{T}
\left\|\frac{\hat a_t^{\mathrm{mot}}-a_t^{\mathrm{mot}}}{s_a}\right\|_2^2.
\end{equation}
Gripper BCE is retained as a diagnostic but excluded from the primary score.
Out-of-fold scoring avoids assigning low difficulty merely because a model
memorized a training trajectory; it measures approximation error for this
policy family, not intrinsic task complexity or failure probability.

\subsection{Curriculum strategies}
Each metric sorts trajectories by ascending score, breaking ties by ID, and
forms 80 Easy, 80 Medium, and 80 Hard trajectories. Forward curricula use
Easy, then Easy+Medium, then all data. Reverse Learner uses Hard, then
Hard+Medium, then all data. Uniform uses all trajectories throughout.
Curriculum boundaries are updates 6,866 and 13,733 of a 20,600-update run.
Within a seed, all methods share the exact initial weights, architecture,
optimizer settings, loss, validation split, and evaluation tasks; only the
eligible trajectory pool differs.

Sampling is uniform over \emph{transitions} in the eligible pool, not uniform
over trajectories. Thus longer trajectories receive more aggregate sampling
probability. The three stages contain 6,866, 6,867, and 6,867 updates. Using
the actual transition counts, cumulative expected group exposures are
59.9/28.1/12.0\% (Task), 59.5/28.3/12.1\% (Expert), and
61.9/27.6/10.6\% (Learner) for Easy/Medium/Hard. Reverse Learner yields
11.6/28.3/60.1\% in that order. Uniform exposure follows each ranking's
transition composition (Task: 31.2/32.9/36.0\%; Expert: 30.6/33.1/36.3\%;
Learner: 34.7/33.6/31.7\%). Consequently, the methods compare staged
pool-sampling schedules that combine ordering with cumulative reweighting; they
do not isolate order under equal per-trajectory exposure.
